\documentclass[11pt]{article}

\usepackage{hyperref}
\usepackage{xurl}
\usepackage{float}
\usepackage{array}
\usepackage[margin=1in]{geometry}
\usepackage[tagged, highstructure]{accessibility}
\usepackage{bookmark}
\usepackage{enumitem}

\hypersetup{
    colorlinks=false,
    pdfborder={0 0 0},
    pdftitle={Sample Threat Model: Protecting a Streamer's Identity},
    pdfauthor={Adrianna Holden-Gouveia},
    pdfsubject={Information Security - Sample threat model for the Encrypted Communications and Threat Models Lab},
    pdflang={en-US},
    pdfkeywords={threat model, sample, risk register, doxxing, account takeover}
}

\bookmarksetup{
    numbered,
    open,
}

\setlist{nosep}

\title{Threat Model: Protecting a Streamer's Identity}
\author{SAMPLE ONLY. Fictional person. Not for submission.}
\date{\vspace{-5ex}}

\begin{document}

\maketitle

\section*{Read This First}
\begin{itemize}
    \item \textbf{This is a sample.} It shows what a threat model can look like in industry. Riley is a made-up person, and the details are invented.
    \item \textbf{Don't copy it.} Copied wording, risks, or controls get no credit. This scenario isn't one of the lab's choices, so you can't use it.
    \item \textbf{Your lab has its own requirements.} Follow the lab instructions, not this layout. Your lab needs EFF's six questions, at least 10 risks, a plan, and at least 5 sources, each with a CRAAP check, a screenshot, and a quote. This sample has no sources section like that.
    \item \textbf{Not professional advice.} If you have a real security problem, get help from someone who can look at your situation.
\end{itemize}

\section*{Accessibility Notice}
This document is also available in HTML format at:

Link: \href{https://aholdengouveia.name/InfoSec/labs/threatmodel-streamer.html}{\textbf{\nolinkurl{https://aholdengouveia.name/InfoSec/labs/threatmodel-streamer.html}}}

\section{Document Control}
\begin{table}[H]
\caption{Document control}
\begin{tabular}{|m{0.3\textwidth}|m{0.6\textwidth}|}
\hline
\textbf{Field} & \textbf{Value} \\
\hline
Document ID & TM-SAMPLE-01 \\
\hline
Version & 1.0 \\
\hline
Date & October 15, 2026 \\
\hline
System owner & Riley (fictional) \\
\hline
Prepared by & Adrianna Holden-Gouveia, for CIS 115 \\
\hline
Classification & SAMPLE. Not for submission. \\
\hline
Next review & April 2027, or sooner if something under When to Review happens
\end{tabular}
\end{table}

\section{Summary}
Riley is a college student who streams video games and posts short videos, with a few thousand followers. The biggest risks are someone taking over Riley's email, and then the channel, and viewers working out Riley's real name and home address, which is called doxxing. The most important controls are a password manager, two-factor authentication with an app, separate accounts for the channel and personal life, and checking photos and videos before posting. After these controls, the remaining risk is Low to Medium, which Riley accepts.

\section{Scope}
\textbf{In scope:} Riley's phone and laptop, their email, streaming, and social media accounts, their messages and calls with friends, family, and sponsors, and anything Riley posts publicly.

\textbf{Out of scope:} the streaming platform's own security, Riley's school accounts, and physical safety at events. These are someone else's responsibility, or need a separate threat model.

\section{Assets}
\begin{table}[H]
\caption{Assets to protect}
\begin{tabular}{|m{0.3\textwidth}|m{0.6\textwidth}|}
\hline
\textbf{Asset} & \textbf{Why it matters} \\
\hline
Real name, home address, and school & If these become public, strangers could show up in person \\
\hline
Personal email account & It can reset the password on almost every other account \\
\hline
Streaming and social media accounts & They hold Riley's followers and income \\
\hline
Private messages & Friends and family share personal details there \\
\hline
Phone and phone number & The phone holds every account, and the number receives some login codes
\end{tabular}
\end{table}

\section{Threat Actors}
\begin{table}[H]
\caption{Who might attack, and why}
\begin{tabular}{|m{0.25\textwidth}|m{0.3\textwidth}|m{0.35\textwidth}|}
\hline
\textbf{Threat actor} & \textbf{What they want} & \textbf{What they can do} \\
\hline
Overly curious or hostile viewers & Riley's real identity and location & Search usernames and photos, send links in chat, ask personal questions \\
\hline
Account takeover scammers & The channel, to resell it or run scams from it & Phishing, password reuse attacks, talking the phone company into moving Riley's number \\
\hline
Fake sponsors & Riley's login, or to install malware & Professional looking emails with ``contracts'' and links \\
\hline
Data brokers & To sell personal data & List names and addresses on people-search websites
\end{tabular}
\end{table}

\section{Assumptions}
\begin{itemize}
    \item The streaming platform and email provider keep their own systems secure.
    \item No one who knows Riley in person is trying to harm them.
    \item Riley's phone and laptop aren't infected with malware today.
    \item Riley isn't a target of a government or a skilled, well-funded attacker.
\end{itemize}
If any of these stop being true, update this threat model.

\section{Entry Points}
These are the ways an attacker could reach Riley's assets:
\begin{itemize}
    \item Login pages for email, streaming, and social media accounts
    \item Riley's posts: photos, videos, the background on stream, and location tags
    \item Usernames shared between personal and channel accounts
    \item Email, direct messages, and links in chat
    \item The phone company's customer service, for moving a phone number
    \item Riley's home internet connection, through its IP address
\end{itemize}

\section{Risk Rating}
Each threat gets a likelihood (how likely it is to happen) and an impact (how bad it would be), from Low to High. The overall risk comes from this table.
\begin{table}[H]
\caption{How likelihood and impact combine into an overall risk}
\begin{tabular}{|m{0.22\textwidth}|m{0.2\textwidth}|m{0.2\textwidth}|m{0.2\textwidth}|}
\hline
\textbf{Likelihood} & \textbf{Low impact} & \textbf{Medium impact} & \textbf{High impact} \\
\hline
High & Medium & High & High \\
\hline
Medium & Low & Medium & High \\
\hline
Low & Low & Low & Medium
\end{tabular}
\end{table}

\section{Threat Register}
Ordered from highest to lowest risk.
\begin{table}[H]
\caption{Threat register}
\begin{tabular}{|m{0.06\textwidth}|m{0.46\textwidth}|m{0.11\textwidth}|m{0.11\textwidth}|m{0.1\textwidth}|}
\hline
\textbf{ID} & \textbf{Threat} & \textbf{Likeli\-hood} & \textbf{Impact} & \textbf{Risk} \\
\hline
T1 & Email takeover, because Riley reuses a password and has no two-factor authentication & High & High & High \\
\hline
T2 & Photos or video reveal Riley's home or school, through location data in the file, or a street sign, school logo, or window view in the background & Medium & High & High \\
\hline
T3 & Location tags or location sharing, such as Snap Map, show where Riley is or goes often & Medium & High & High \\
\hline
T4 & The same username on personal and channel accounts links the channel to Riley's real name & High & Medium & High \\
\hline
T5 & A fake sponsor email delivers malware, or a fake login page steals the channel password & Medium & High & High \\
\hline
T6 & An attacker talks the phone company into moving Riley's number to a new SIM card, and gets text message login codes & Low & High & Medium \\
\hline
T7 & A link in chat records Riley's IP address, which shows their city and internet provider & Medium & Medium & Medium \\
\hline
T8 & People-search websites list the family's address under Riley's real name & Medium & Medium & Medium \\
\hline
T9 & Friends tag Riley's real name in comments or posts & Medium & Medium & Medium \\
\hline
T10 & Phone apps have permissions they don't need, such as location or contacts, and that data is leaked or sold & Medium & Low & Low
\end{tabular}
\end{table}

\section{Controls}
\begin{table}[H]
\caption{Controls, and the threats each one reduces}
\begin{tabular}{|m{0.06\textwidth}|m{0.6\textwidth}|m{0.12\textwidth}|m{0.1\textwidth}|}
\hline
\textbf{ID} & \textbf{Control} & \textbf{Threats} & \textbf{Priority} \\
\hline
C1 & Use a password manager, with a different strong password for every account, starting with email & T1, T5 & Now \\
\hline
C2 & Turn on two-factor authentication with an authenticator app, not text messages & T1, T5, T6 & Now \\
\hline
C3 & Make a separate email address for the channel and sponsors, and a different username for personal accounts. Make personal accounts private. & T1, T4, T9 & Now \\
\hline
C4 & Before posting, check the background for signs, mail, logos, and window views. Turn off location for the camera app, and stop tagging locations. & T2, T3 & Now \\
\hline
C5 & Turn off Snap Map and other location sharing. Post about places after leaving them. & T3 & Now \\
\hline
C6 & Never open unexpected files. Check that a sponsor is real by finding the company's website yourself, not through the email's link. & T5 & Now \\
\hline
C7 & Add a PIN or passcode to the phone account, so the phone company won't move the number without it & T6 & This month \\
\hline
C8 & Use a well-known, paid VPN that has had an independent audit of its no-logs promise, especially when opening links from chat or playing online games & T7 & This month \\
\hline
C9 & Ask people-search websites to remove the family's listing, and ask friends not to tag Riley's real name & T8, T9 & This month \\
\hline
C10 & Review app permissions every few months, and turn off what isn't needed & T10 & Ongoing
\end{tabular}
\end{table}

\subsection*{Communications}
\textbf{Messages with friends and family:} Signal, which is end-to-end encrypted, so only the people in the chat can read the messages. \textbf{Sponsors:} the channel email only. \textbf{Phone calls and texts:} regular calls and texts aren't end-to-end encrypted, so Riley gives sponsors a free internet phone number, such as Google Voice, and keeps their real number for people they know.

\subsection*{VPN decision}
A VPN is worth it here, for one reason: it hides Riley's home IP address from websites and games (T7). It doesn't protect against stolen passwords, phishing, or a photo that shows Riley's street, so it's a lower priority than C1 to C6.

\section{Accepted Risk}
After the controls, every threat is rated Low or Medium. Riley accepts these risks:
\begin{itemize}
    \item Someone could still guess Riley's city from things they mention on stream, such as the weather.
    \item A platform data breach could leak Riley's channel email. This is why it's separate from their personal email.
\end{itemize}
Riley decided against Tor and a second phone. They'd be more trouble than they're worth for this threat level.

\section{Not Covered}
This model doesn't cover a viewer who finds Riley in person, or threats of harm. If that happens, Riley will report it to the platform and the police, and update this model.

\section{When to Review}
Review this threat model every 6 months, and right away if:
\begin{itemize}
    \item The channel grows a lot, or Riley starts showing their face or real name
    \item An account is taken over, or Riley's address is posted online
    \item Riley moves, changes schools, or gets a new phone
\end{itemize}

\section{References}
EFF Surveillance Self-Defense: Your Security Plan
Link: \href{https://ssd.eff.org/module/your-security-plan}{\textbf{\nolinkurl{https://ssd.eff.org/module/your-security-plan}}}

EFF Surveillance Self-Defense: Protecting Yourself on Social Networks
Link: \href{https://ssd.eff.org/module/protecting-yourself-social-networks}{\textbf{\nolinkurl{https://ssd.eff.org/module/protecting-yourself-social-networks}}}

The lab this sample goes with
Link: \href{https://www.aholdengouveia.name/InfoSec/labs/enccomm.html}{\textbf{\nolinkurl{https://www.aholdengouveia.name/InfoSec/labs/enccomm.html}}}

\end{document}
